\section{Spherical consequences}\label{sec:consequences}

The global sphere estimate has now been proved. We give the complete
transfer to the translated-tube maximal operator defined in the
introduction, retaining the weighted tube estimate needed for its
strong operator norm. This uses sphere factorization in addition to
Theorem~\ref{thm:main}.

\subsection{The Kakeya maximal transfer}

\begin{proof}[Proof of Corollary~\ref{cor:sphere-kakeya-maximal}]
Sphere factorization, due to Bourgain~\cite{Bourgain1991Besicovitch},
upgrades Theorem~\ref{thm:main} to
\(\|Eg\|_{L^q(\R^3)}\lesssim_q\|g\|_{L^q(S^2)}\) for every \(q>3\);
see~\cite[Section~1.1, (1.1)--(1.2)]{Buschenhenke2024} for the
weak-type factorization and interpolation in the sphere case.
Fix \(3<q<4\) and put $r=q/2$. Consider unit tubes
$T_j=T_\delta(x_j,\omega_j)$ in $\delta$-separated directions.
Choose smooth nonnegative bumps $b_j$, bounded by $1$, supported on
disjoint caps of radius $c\delta$ about $\omega_j$, and equal to $1$
on the concentric caps of radius $c\delta/2$. Here $c>0$ is a fixed
small constant. The modulation
\[
 h_j(\zeta)=e^{-2\pi i\delta^{-2}x_j\cdot\zeta}b_j(\zeta)
\]
translates the extension to the center of $\delta^{-2}T_j$.
We have $\|h_j\|_q^q\lesssim\delta^2$ and
$|Eh_j|\gtrsim\delta^2$ on that dilated tube. To see the lower bound,
write $y-\delta^{-2}x_j=s\omega_j+v$, where
$|s|\le\delta^{-2}/2$, $v\perp\omega_j$, and $|v|\le\delta^{-1}$.
On the cap, the phase after removing its value at $\omega_j$ has
magnitude at most $C(c+c^2)$. Taking $c$ sufficiently small makes
the real part of the remaining integrand a fixed positive fraction
of $b_j$, whose integral is comparable to $\delta^2$.

For arbitrary weights $a_j\ge0$ and independent random signs
$\varepsilon_j$, Khintchine's inequality and the diagonal estimate give
\[
 \delta^{2q-6}\Big\|\sum_j a_j1_{T_j}\Big\|_r^r
 \lesssim_q
 \mathbb E\Big\|E\Big(\sum_j\varepsilon_j a_j^{1/2}h_j\Big)\Big\|_q^q
 \lesssim_q\delta^2\sum_j a_j^r.
\]
The factor $\delta^{-6}$ comes from $y=\delta^{-2}x$, while the
last inequality uses the disjoint cap supports. Taking the $r$th
root gives the weighted tube estimate
\[
 \Big\|\sum_j a_j1_{T_j}\Big\|_{L^r(\R^3)}
 \lesssim_q \delta^{12/q-4}
       \Big(\delta^2\sum_j a_j^r\Big)^{1/r}.
\]

Set $p=r'=q/(q-2)$. For nonnegative $f\in L^p(\R^3)$ and $c_j\ge0$,
the normalization $|T_j|=\pi\delta^2$ gives
\[
 \delta^2\sum_j c_j\frac1{|T_j|}\int_{T_j}f
 =\frac1\pi\int f\sum_jc_j1_{T_j}
 \lesssim_q\delta^{12/q-4}\|f\|_p
                \Big(\delta^2\sum_jc_j^r\Big)^{1/r}.
\]
Duality for the discrete measure assigning mass $\delta^2$ to each
direction bounds the $p$-norm of these tube averages. The bound is
uniform in all tube locations, so each $T_j$ may independently
approximate the supremum defining $K_\delta f(\omega_j)$.

Choose a maximal $\delta$-separated net on $S^2$, with angular cells
of area comparable to $\delta^2$. Every direction in a cell is within
$\delta$ of its center $\omega_j$. For $\delta<1/4$, the corresponding
tubes satisfy
\[
 T_\delta(a,\omega)\subset
 T_{2\delta}(a-\omega_j/4,\omega_j)
 \cup T_{2\delta}(a+\omega_j/4,\omega_j).
\]
Indeed, the transverse displacement is at most $3\delta/2$ and the
axial displacement at most $1/2+\delta^2<3/4$. Thus
$K_\delta f(\omega)\le8K_{2\delta}f(\omega_j)$.
Sphere packing splits the net into a fixed number of
$2\delta$-separated families. Apply the discrete bound at width
$2\delta$ to each family and sum over the angular cells. This gives
Bourgain's strong maximal transfer, in the form stated in
\cite[Theorem~3.3, pp.~10--11]{MattilaFourierNotes}:
\[
 \|K_\delta f\|_{L^p(S^2)}
 \lesssim_q\delta^{12/q-4}\|f\|_{L^p(\R^3)},
 \qquad p=\frac{q}{q-2}.
\]
For $\delta\ge1/4$, the same bound follows directly from H\"older's
inequality. Interpolating with
\(\|K_\delta f\|_\infty\le\|f\|_\infty\), at parameter \(p/3\), yields
\[
 \|K_\delta f\|_{L^3(S^2)}
 \lesssim_q\delta^{-4(q-3)/(3(q-2))}\|f\|_{L^3(\R^3)}.
\]
Choosing \(q>3\) sufficiently close to \(3\) proves the claim.
\end{proof}

This deduction starts from the spherical conclusion of
Theorem~\ref{thm:main} and uses sphere-specific factorization; it does
not invoke the general-surface diagonal companion or the independent
maximal theorem.


\subsection{The open mixed-norm range on the sphere}

The diagonal sphere estimates also give the strict mixed-norm range.
For $1<a\le\infty$ and $3<b<\infty$, write $a'=a/(a-1)$ when
$a<\infty$ and $a'=1$ when $a=\infty$. Then
\begin{equation}\label{eq:sphere-mixed-norm}
 \|Eg\|_{L^b(\R^3)}\le C_{a,b}\|g\|_{L^a(S^2)},
 \qquad b>2a',
\end{equation}
for every complex $g\in L^a(S^2)$. The defining integral for $E$
remains meaningful because $\sigma(S^2)<\infty$ implies
$L^a(S^2)\subset L^1(S^2)$. Thus the conclusion covers a larger
class of input densities, with an output exponent depending on
their integrability.

Indeed, for $1<a\le b$, set $r=1+b/a'>3$ and $\vartheta=r/b\le1$.
Interpolate the diagonal $L^r(S^2)\to L^r(\R^3)$ estimate from
sphere factorization with the elementary $L^1(S^2)\to L^\infty(\R^3)$
estimate. The identities
\[
 \frac1b=\frac{\vartheta}{r},\qquad
 \frac1a=1-\vartheta+\frac{\vartheta}{r}
\]
give \eqref{eq:sphere-mixed-norm}; when $a=b$, the diagonal estimate
itself suffices. For $b<a<\infty$, use the diagonal estimate at $b$
and H\"older's inequality on $S^2$. The case $a=\infty$ is
Theorem~\ref{thm:main}. This argument uses the strict inequality
$b>2a'$ to choose $r>3$ and makes no assertion on the scaling line
$b=2a'$.
